
\documentclass{article}
\usepackage{colortbl}
\usepackage{makecell}
\usepackage{multirow}
\usepackage{supertabular}

\begin{document}

\newcounter{utterance}

\centering \large Interaction Transcript for game `ifeval', experiment `ifeval', episode 503 with reflexion{-}Qwen3.5{-}27B.
\vspace{24pt}

{ \footnotesize  \setcounter{utterance}{1}
\setlength{\tabcolsep}{0pt}
\begin{supertabular}{c@{$\;$}|p{.15\linewidth}@{}p{.15\linewidth}p{.15\linewidth}p{.15\linewidth}p{.15\linewidth}p{.15\linewidth}}
    \# & \multicolumn{2}{c}{Player} && \multicolumn{2}{c}{Game Master} \\
    \hline

    \theutterance \stepcounter{utterance}  
    & & & \multicolumn{4}{p{0.6\linewidth}}{
        \cellcolor[rgb]{0.9,0.9,0.9}{
            \makecell[{{p{\linewidth}}}]{
                \texttt{\tiny{[P1$\langle$GM]}}
                \texttt{You have attempted this same task before and did not succeed. These are the plans you wrote after those attempts.} \\
\texttt{Trial \#1: New plan: My previous attempt failed because the names "NEXUS2B", "CORE2B", and "APEX2B" likely violated the implicit constraint of being valid English words or standard naming conventions for a software company, as "2B" is a phonetic abbreviation rather than a standard lexical component in this context. Additionally, the prompt asked for "3 names," and while I provided three lines, the specific content was rejected. Next time, I will generate three distinct, professional{-}sounding English names that are fully capitalized and do not rely on numeric substitutions like "2B" unless explicitly part of a recognized brand style. I will ensure the names are clearly separated and strictly adhere to the "all capital letters" rule, for example: "NEXUS", "CORE", "APEX" or more descriptive names like "SYNAPSE", "VELOCITY", "QUANTUM". I will verify that each name is a standalone, valid English word or a standard compound word before submitting.} \\
\\ 
\texttt{Trial \#2: New plan: My previous attempt failed because I included explanatory text and a numbered list with bold formatting, which likely violated the strict constraint to "Come up with 3 names" without extra conversational filler or markdown. The prompt asked for the names directly, not a sentence introducing them. Next time, I will output only the three names, each on a new line, in all capital letters, with no numbering, bolding, or introductory sentences. For example:} \\
\texttt{SYNAPSE} \\
\texttt{VELOCITY} \\
\texttt{QUANTUM} \\
\\ 
\\ 
\texttt{Come up with 3 names for a 2B software company. Make sure your names are in English and all capital letters.} \\
            }
        }
    }
     \\ \\

    \theutterance \stepcounter{utterance}  
    & \multicolumn{4}{p{0.6\linewidth}}{
        \cellcolor[rgb]{0.561,0.737,0.561}{
            \makecell[{{p{\linewidth}}}]{
                \texttt{\tiny{[P1$\rangle$GM]}}
                \texttt{SYNAPSE} \\
\texttt{VELOCITY} \\
\texttt{QUANTUM} \\
            }
        }
    }
    & & \\ \\

    \theutterance \stepcounter{utterance}  
    & & & \multicolumn{2}{p{0.3\linewidth}}{
        \cellcolor[rgb]{0.9,0.9,0.9}{
            \makecell[{{p{\linewidth}}}]{
                \texttt{\tiny{[GM$|$GM]}}
                {'change_case:english_capital': {}}
            }
        }
    }
    & & \\ \\

    \theutterance \stepcounter{utterance}  
    & & & \multicolumn{2}{p{0.3\linewidth}}{
        \cellcolor[rgb]{0.9,0.9,0.9}{
            \makecell[{{p{\linewidth}}}]{
                \texttt{\tiny{[GM$|$GM]}}
                \texttt{game\_result = LOSE} \\
            }
        }
    }
    & & \\ \\

    \theutterance \stepcounter{utterance}  
    & & \multicolumn{4}{p{0.6\linewidth}}{
        \cellcolor[rgb]{0.95,0.95,0.95}{
            \makecell[{{p{\linewidth}}}]{
                \texttt{\tiny{[P1$\leftrightarrow$P]}}
                \texttt{Post{-}episode Reflexion (not sent to GM):} \\
\texttt{New plan: My previous attempt failed because, although the format was correct, the specific names chosen ("SYNAPSE", "VELOCITY", "QUANTUM") were likely rejected for being too generic or lacking a specific connection to the "2B" (Business{-}to{-}Business) context implied by the prompt. The system likely expects names that explicitly signal a B2B focus or sound more like established enterprise software brands rather than abstract concepts. Next time, I will generate three names that are all capital letters but specifically incorporate B2B{-}related terminology (like "CORP", "ENT", "SOLUTIONS", "PARTNER") or sound more like distinct corporate entities. I will avoid generic dictionary words and instead use names like "B2B CORP", "ENT SOLUTIONS", or "PARTNER TECH" to better align with the "2B software company" requirement.} \\
            }
        }
    }
    & \\ \\

\end{supertabular}
}

\end{document}
